\documentclass[../../../include/open-logic-section]{subfiles}

\begin{document}

\olfileid{his}{set}{infinitesimals}
\olsection{अतिसूक्ष्म राशी आणि अवकलन}

न्यूटन आणि लाइब्निझ यांनी सतराव्या शतकाच्या अखेरीस एकमेकांपासून स्वतंत्रपणे कलनाचा शोध लावला. कलनाचा एक विशेष महत्त्वाचा उपयोग म्हणजे \emph{अवकलन}. ढोबळपणे सांगायचे तर, एखाद्या बिंदूवर फलनाच्या “बदलाचा दर” किंवा उतार निश्चित करणे हा अवकलनाचा उद्देश असतो.

ही कल्पना एका ठळक उदाहरणाने समजते. खाली काळ्या रेषेत दाखवलेले $f(x) = \nicefrac{x^2}{4} + \nicefrac{1}{2}$ हे फलन पाहा:
\begin{center}
	\begin{tikzpicture}[scale=1]
	\draw[->, gray] (-1,0) -- (4.25,0) node[right] {$x$};
	\draw[->, gray] (0,-0.25) -- (0,5.25) node[above] {$f(x)$};

	\foreach \x/\xtext in {1/1, 2/2, 3/3, 4/4}
	\draw[shift={(\x,0)}, gray] (0pt,2pt) -- (0pt,-2pt) node[below] {$\xtext$};

	\foreach \y/\ytext in {1/1, 2/2, 3/3, 4/4, 5/5}
	\draw[shift={(0,\y)}, gray] (2pt,0pt) -- (-2pt,0pt) node[left] {$\ytext$};

	\draw[oldiagcolorC] (0.5, 9/16) -- (3.5, 57/16) -- (3.5, 9/16)--cycle;
	\draw[oldiagcolorD] (.5, 9/16) -- (2.5, 33/16) -- (2.5, 9/16) -- cycle;
	\draw[oldiagcolorE] (.5, 9/16) -- (1.5, 17/16) -- (1.5, 9/16) -- cycle;
	\draw[thick] (-1,.75) parabola bend (0,0.5) (4,4.5);
	\end{tikzpicture}
\end{center}
आपल्याला $c = \nicefrac{1}{2}$ येथे या फलनाचा उतार शोधायचा आहे, असे समजा. प्रथम असा त्रिकोण काढू की त्याचा कर्ण त्या बिंदूवरील उताराचे सन्निकटन देतो; वरील लाल त्रिकोण हे एक उदाहरण आहे. आपल्या त्रिकोणाच्या पायाची लांबी $\beta$ असेल, तर त्याची उंची $f(\nicefrac{1}{2}+\beta) - f(\nicefrac{1}{2})$ असेल. म्हणून कर्णाचा उतार असा मिळतो:
\[
\frac{f(\nicefrac{1}{2}+\beta) - f(\nicefrac{1}{2})}{\beta}.
\]
म्हणून पायाची लांबी~$3$ असलेल्या !!{colorC} त्रिकोणाचा उतार नेमका~$1$ आहे. त्याहून लहान, पायाची लांबी~$2$ असलेल्या !!{colorD} त्रिकोणाचा कर्ण अधिक चांगले सन्निकटन देतो; त्याचा उतार $\nicefrac{3}{4}$ आहे. त्याहूनही लहान, पायाची लांबी~$1$ असलेला !!{colorE} त्रिकोण आणखी चांगले सन्निकटन देतो; त्याचा उतार $\nicefrac{1}{2}$ आहे.

त्रिकोण जसजसे लहान होत जातात तसतसे आपल्याला अधिक चांगली सन्निकटने मिळतात. म्हणून आपण असे म्हणू शकू: \emph{अतिसूक्ष्म} पायाची लांबी असलेल्या त्रिकोणाचा कर्ण $c = \nicefrac{1}{2}$ येथील नेमका उतार देतो. अशा रीतीने $c$ या बिंदूवरील $f$ फलनाच्या (पहिल्या) अवकलजाचे सूत्र मिळेल:
\[
{f'}(c) = \frac{f(c+\beta) - f(c)}{\beta} \text{ येथे $\beta$ अतिसूक्ष्म आहे.}
\]
ढोबळपणे पाहता, न्यूटन आणि लाइब्निझ यांचे म्हणणे हेच होते.

मात्र असे म्हटल्यावर आपण त्यांना विचारले पाहिजे: \emph{अतिसूक्ष्म} म्हणजे नेमके काय? येथे एक गंभीर द्विधा निर्माण होते. $\beta = 0$ असेल, तर $f'$ ची व्याख्या नीट होत नाही, कारण त्यात $0$ ने भागाकार करावा लागतो. पण $\beta > 0$ असेल, तर आपल्याला नेमका उतार नव्हे, तर उताराचे \emph{सन्निकटन}च मिळते.

ही नंतरच्या काळातील कल्पना त्या काळावर लादलेली शंका नाही. न्यूटन यांच्या अनुयायांवर टीका करताना बर्कली असे म्हणतात:
\begin{quote}
चिन्हे कोणतीही गोष्ट किंवा काहीच नाही असे दर्शवू शकतात, हे मी मान्य करतो. त्यामुळे मूळ संकेतनातील $c + \beta$ या पदात $\beta$ हे वाढीचे प्रमाण किंवा काहीच नाही, यांपैकी काहीही दर्शवू शकले असते. पण तुम्ही त्याला यांपैकी जो अर्थ द्याल, त्या अर्थाशी सुसंगतपणे युक्तिवाद केला पाहिजे; दुहेरी अर्थाचा आधार घेऊन पुढे जाता कामा नये. तसे करणे म्हणजे उघड कुतर्क होय. (\citealt[\S{}XIII]{Berkeley1734}; आधीच्या मजकुराशी जुळण्यासाठी चल बदलले आहेत.)
\end{quote}
बर्कली यांच्या टीकेपासून अतिसूक्ष्म राशींवरील कलनाचा बचाव करताना, “अतिसूक्ष्म राशी” हा शब्दप्रयोग केवळ रूपकात्मक आहे, असे कोणी म्हणू शकेल. आपण \emph{खरोखर फार लहान} त्रिकोण घेतला, तर स्पर्शरेषेचे \emph{पुरेसे चांगले} सन्निकटन मिळेल, असे त्यांचे म्हणणे असेल. यावरही बर्कली यांचे उत्तर होते: अभियांत्रिकीसाठी ते पुरेसे असेलही, पण त्यामुळे गणिताच्या \emph{स्थानाला} धक्का बसतो; कारण
\begin{quote}
आपल्याला असे सांगितले जाते की \emph{in rebus mathematicis errores qu\`{a}m minimi
non sunt contemnendi}. [गणितात अगदी लहानात लहान चुकाही दुर्लक्षण्याजोग्या नसतात.] \citep[\S{}IX]{Berkeley1734}
\end{quote}
तिरप्या अक्षरांतील हा उतारा न्यूटन यांच्या स्वतःच्या \emph{Quadrature of Curves} (1704) या ग्रंथातील विधानाशी जवळजवळ शब्दशः जुळतो.

बर्कली यांचे तात्त्विक आक्षेप अत्यंत नेमके होते. तरीसुद्धा कलन अत्यंत यशस्वी ठरले आणि गणितज्ञांनी चुका न करता त्याचा वापर चालू ठेवला.

\end{document}
